\scanpage{343}
% top centre
\mspage{259}
% line 1
On a donc prouvé l'essentiel du

% lines 2–11
\uline{Lemme \unsure{4|2}.10.7} Sous les conditions de 4.10.3 a) b), et pour un système $X(\cdot)\in\operatorname{Ob}\prod_{j\in\operatorname{Ob}J}\unsure{\mathcal E|\mathscr E|\mathcal C}_j$ donné, on peut trouver un système inductif $\bigl(X_\alpha(\cdot)\bigr)_\alpha$ dans $\prod\mathcal E_j$, indexé par les ordinaux $\alpha$ tels que $\mathrm{card}\,\alpha\leq\pi$, \deleted{\unsure{et pour}}\inserted{et} pour toute flèche $\lambda:j'\unsure{\mathrel{-}|\longrightarrow}j$ dans $J$ \deleted{\unsure{une}} des flèches $t_\alpha(\lambda):X_\alpha(j')\unsure{\mathrel{-}|\longrightarrow}X_\alpha(j)$, de telle façon que l'on ait les conditions suivantes\unsure{:|.}

% line 12
\[
 \uline{1}^{\circ})\quad X_0(\cdot)=X(\cdot)
\]

% line 13
$2^\circ)$ Les conditions (i) (ii) (iii) (p. 256)

% lines 14–17
$3^\circ)$ Si on se donne $F$ dans $\operatorname{Ob}\underline{\unsure{\Gamma|\mathscr T}}(\mathcal E/J)$, et des flèches $\unsure{u_j|u(j)}:F(j)\longrightarrow X(j)\underset{\unsure{\lambda|{}}}{=}X_0(j)$, \unsure{sps|tqs} que pour \unsure{tte} flèche $\lambda$\inserted{$:j'\longrightarrow j$} dans $J$, on a
% bottom, label beside diagram
\diagram{\begin{tikzcd}[ampersand replacement=\&,column sep=1.15em,row sep=3ex]
 {} \arrow[d,phantom,"{\text{\unsure{(iv)|(4)}}}" description] \& F(j') \arrow[r,"{u(j')}"] \arrow[d,"{F(\lambda)}"'] \& X_0(j') \arrow[dr,"{t_0(\lambda)}"] \& {} \& {} \arrow[d,phantom,"{\text{\uline{commutatif}}}" description] \\
 {} \& F(j) \arrow[r,"{u(j)}"] \& X_0(j) \arrow[r,"{\tau_0(j)}"] \& X_1(j\unsure{{}^{\prime\prime}|{}}) \& \cdot\hspace{2.6em}\cdot
\end{tikzcd}}
% archive sheet number: 342
